%--------------------------------------------------------------------%
% BAB IV IMPLEMENTASI PROTOTIPE
%--------------------------------------------------------------------%
\chapter{IMPLEMENTASI PROTOTIPE} \label{bab:4}
\petunjuk{Wujud solusi yang dibangun. Pada seminar proposal boleh berupa prototipe awal bila ada. Uraikan rinci komponen yang Anda kerjakan; komponen anggota lain cukup diringkas. Bagian ini wajib ditulis sendiri dan berbeda dari laporan anggota tim lain (Pedoman TA 5.1.3). \tahap{\opsional{Opsional (prototipe awal)}}{\wajib{WAJIB (final)}}}

\section{Arsitektur dan Komponen} \label{sec:4.1}
\petunjuk{Arsitektur yang diimplementasikan dan komponen utamanya.}
\tulisuraian

\section{Fitur dan Antarmuka} \label{sec:4.2}
\petunjuk{Fitur yang tersedia, dilengkapi tangkapan layar atau foto prototipe sebagai gambar bernomor.}
\tulisuraian

\section{Konfigurasi dan Reproduktibilitas} \label{sec:4.3}
\petunjuk{Uraikan perangkat keras, perangkat lunak dan versinya, repositori kode dan dokumentasi teknis, serta cara menjalankan atau memasang agar prototipe dapat dibangun ulang.}
\tulisuraian
